\documentclass[11pt,a4paper]{engineering_report}

\institucion{Departamento de Ingenieria Mecanica}
\facultad{Area de Diseno Mecanico y Simulacion CAE}
\asignatura{Elementos de Maquinas / Resistencia de Materiales}
\codigo{MEM-CAE-01}
\titulo{Analisis Estructural y Dimensionamiento de Transmision por Engranes}
\autor{Estudiante / Integrantes del Proyecto}
\periodo{Periodo Academico 2026}
\logotipo{logo_studio.png}

\begin{document}
\makeengineeringcover

\section{Resumen Ejecutivo}
Este informe presenta la memoria de calculo analitico y la verificacion dimensional de una transmision mecanica bajo normativa ISO 53, DIN 6885 e ISO 281.

\begin{sciformula}{Ecuacion de Vida Nominal de Rodamientos (ISO 281)}
\[
L_{10h} = \frac{10^6}{60\,n} \left( \frac{C}{P_{\mathrm{eq}}} \right)^p, \qquad p = 3 \text{ (bolas)}
\]
\end{sciformula}

\section{Tabla de Resultados Normalizados}
\begin{table}[H]
\centering
\caption{Resumen dimensional y de vida util verificada.}
\begin{tabular}{@{}lccc@{}}
\toprule
\textbf{Componente} & \textbf{Norma} & \textbf{Parametro Critico} & \textbf{Valor} \\
\midrule
Pinon recto $z_1=20$ & ISO 53 & Diametro primitivo $d_1$ & $40.00\text{ mm}$ \\
Ajuste eje-buje & ISO 286 & Tolerancia $\mathrm{H7/h6}$ & $+25 / -13\,\mu\text{m}$ \\
Rodamiento 6205 & ISO 281 & Vida nominal $L_{10h}$ & $> 25{,}000\text{ h}$ \\
\bottomrule
\end{tabular}
\end{table}

\begin{scibox}{Conclusion Tecnica}
Todos los coeficientes de seguridad superan el umbral minimo requerido ($n_s \ge 2.0$), manteniendo compatibilidad directa con la manufactura convencional.
\end{scibox}

\end{document}
